\section{Introduction}\label{sec:introduction}
At the nominal false-alarm level $\alpha=0.20$, the target-only alarm studied in this paper reaches an empirical FAR of 0.341 on Gaussian-corrupted MVTec at $k=4$. This is the largest dataset-level aggregate in the frozen corruption grid and is approximately 1.7 times the requested level. It is not a population-level failure probability: it is a benchmark observation pooled over categories and five support-sampling seeds under one declared synthetic shift. Nevertheless, it exposes an operational gap that ranking metrics alone cannot reveal. A detector may assign defects higher anomaly scores than normal images and obtain a strong area under the receiver operating characteristic curve (AUROC) or average precision (AP), while a fixed score threshold still produces substantially different false-alarm behavior across categories or input conditions.

This gap is particularly consequential in few-shot industrial anomaly detection (AD). A target category commonly provides only $k\in\{1,2,4,8\}$ known-normal support images, whereas anomalous examples are rare, heterogeneous, and unavailable for threshold fitting. Frozen visual foundation models make ranking feasible in this regime: DINOv2~\cite{oquab2024dinov2} patch features support nearest-neighbor (NN) memory banks, as in PatchCore-style methods and AnomalyDINO~\cite{roth2022towards,damm2025anomalydino}, or compact PCA normal subspaces, as in SubspaceAD~\cite{subspacead2026}. These advances do not by themselves specify which alarm threshold has a stable operational meaning.

Two responses are natural. The first is to calibrate on the $k$ target-normal supports. Our target-only route evaluates each support image in a LOIO fit and ranks a test score against the resulting normal scores. This construction is label-free and interpretable, but its rank values occupy only the grid $\{j/(k+1)\}$. Increasing $k$ refines that grid; it does not ensure that support and shifted test scores satisfy the exchangeability conditions needed for false-alarm validity. The empirical FAR of 0.341 is an example of this distinction: the operating level is attainable, yet the shifted-normal rank distribution is anti-conservative relative to the exchangeable reference.

The second response is to borrow normal evidence from non-target source categories. A larger source archive offers finer empirical ranks and can propose thresholds below the target-only grid. It also changes the statistical question. Independent image draws from a declared source mixture can estimate risk for that mixture, whereas a guarantee intended for a previously unseen category depends on iid category units from a declared meta-population. Repeated images, corruption views, or support seeds do not create new category units. The central question of this paper is therefore:

\begin{quote}
How much source evidence, measured in which statistical unit, is needed to certify an anomaly threshold for a new category?
\end{quote}

We answer this question at three levels. First, for any deterministic, uniformly valid, distribution-free 95\% UCB based only on iid category losses, an all-zero sample cannot certify $\alpha=0.20$, 0.10, or 0.05 with fewer than 14, 29, or 59 categories, respectively, even without multiplicity. These are optimistic impossibility-side limits, not sufficient sample sizes for the implemented procedure. Under the frozen simultaneous allocation, the corresponding lower limits are at least 22, 46, and 94 for one candidate, while the declared Hoeffding rule requires 60, 240, and 958. Second, we audit target-only normal rank values under four controlled corruptions and show that crossing the finite resolution floor does not imply a stable FAR\@. Third, we use CRESS to separate source categories into disjoint reference, threshold-proposal, and certification roles and test whether the available archive reaches the predicted feasibility boundary.

CRESS is a source-assisted thresholding and audit protocol, not a new anomaly ranker. The inherited ranking substrate is a frozen DINOv2 PCA residual. The reference categories define a source-reference rank map, the proposal categories define a finite candidate set, and the certification categories evaluate fixed candidates with a family-adjusted UCB\@. The protocol returns the largest passing threshold or the fail-closed value zero. Source-risk dominance can transfer the bound to a particular target category. An iid category-sampling model supports only a marginal bound averaged over a new-category draw, not category-conditional control for every realized target. Normalization and source pooling establish neither condition.

The strict audit contains only three certification categories in VisA cells and four in MVTec and transfer cells. This shortage is not solely a consequence of the frozen $0.50/0.25/0.25$ split. The eligible source pools contain 11, 14, or 15 categories, so keeping the reference and proposal roles nonempty leaves at most 9, 12, or 13 certification categories, all below the optimistic 14-category limit at $\alpha=0.20$. Accordingly, category-unit certification assigns $\tau^\star=0$ to every target cell across the 960 frozen configurations. This is not 960 independent failures, nor a compute or optimization failure. The configurations span jobs, source-view modes, support sizes, nominal levels, and conditions, but share the same category-count obstruction. Their value is to show that the fail-closed outcome persists throughout the declared grid in the regime where the analytic calculation predicts non-feasibility. By contrast, under the idealized iid source-image model, image-unit analyses have 72 to 191 units per cell and select positive thresholds in 36.7\% to 60.3\% of target cells, depending on the source-to-target job. This contrast is not a stronger-versus-weaker version of one certificate: the two analyses target different risks. The same source archive is statistically large when counted in images but insufficient when the deployment claim concerns marginal risk for a new-category draw.

The paper makes three contributions:
\begin{itemize}
    \item \textbf{A category-count feasibility calculus.} We derive the Hoeffding-specific requirement and a broader all-zero lower bound for deterministic uniformly valid distribution-free category-risk upper bounds. The resulting calculations quantify the minimum budget of iid category draws implied by the requested risk and confidence levels and serve as a study-design tool rather than a post hoc explanation alone.
    \item \textbf{An operational shift audit.} Across 15 MVTec and 12 VisA categories, the target-only analysis distinguishes unattainable nominal levels from attainable but unstable operating points. The largest dataset-level aggregate FAR is 0.341 on Gaussian-corrupted MVTec at $k=4$ and $\alpha=0.20$; VisA is conservative in the corresponding cells, showing that the shift effect is dataset-dependent rather than universal.
    \item \textbf{An estimand-aware source protocol.} CRESS separates reference, proposal, and certification categories, retains zero fallbacks, and reports image-mixture and category-level conclusions without substituting one for the other. The frozen strict audit verifies the predicted category shortage within datasets and for transfer from MVTec to VisA and the Metal Parts Defect Dataset (MPDD); direct pooling and compact-ranker results are reported only as supporting empirical context.
\end{itemize}

Figure~\ref{fig:certificate-feasibility} previews the central design calculation. The black curve is the optimistic multiplicity-free floor for deterministic uniformly valid distribution-free bounds in the scope of Proposition~3; the blue and orange curves add the frozen family allocation and the declared Hoeffding rule in their most favorable one-candidate case. A positive certificate is infeasible within the corresponding rule class whenever its curve remains above the requested $\alpha$. The shaded region marks the three or four certification categories available in the strict audit; Table~\ref{tab:certificate-feasibility} later reports the exact candidate-family calculations.

\begin{figure*}[pos=t]
    \centering
    \includegraphics[width=0.97\linewidth]{fig_certificate_feasibility}
    \caption{Analytic category-count feasibility after observing zero loss in every certification category. The black curve is the multiplicity-free lower bound $1-0.05^{1/n}$ for deterministic uniformly valid distribution-free procedures in Proposition~3. The blue curve applies the frozen allocation $\beta=\delta/(2AM)$ with $\delta=0.05$, $A=3$ reported levels, and the favorable one-candidate case $M=1$; the orange curve is the corresponding Hoeffding floor. Horizontal lines mark the requested risk levels, and the shaded region marks the three or four distinct certification categories available in the frozen audit; interpreting them as iid units is a modeling assumption. The black circular crossings at the 14, 29, and 59 axis ticks mark optimistic necessary counts without multiplicity, not sufficient sample sizes for CRESS\@. A curve above a requested $\alpha$ rules out a positive threshold only within its stated procedure class, even under all-zero observed category loss.}
    \label{fig:certificate-feasibility}
\end{figure*}

The remainder of the paper is organized as follows. Section~\ref{sec:related} positions the study against few-shot detectors, calibration, and conformal risk control. Section~\ref{sec:method} defines the inherited scorer, target-only route, source construction, CRESS protocol, and certificate scope. Section~\ref{sec:experiments} specifies the frozen datasets, statistical units, partitions, and claim rules. Section~\ref{sec:results} proceeds from the shift audit to category-count feasibility and strict CRESS, followed by the supporting ranking and storage analysis. Sections~\ref{sec:limitations} and~\ref{sec:conclusion} state the boundaries and implications.
